\chapter{Discussion \& Conclusion}\label{chap:conclusion}

In this thesis, we introduced our software package, \textbf{TriLLIEM}, for fitting the Weinberg et al.~\cite{Wein+98} log-linear model, compared it to existing software, and used \textbf{TriLLIEM} to assess the log-linear model in circumstances that have not been thoroughly studied before.
We designed \textbf{TriLLIEM} to fill the gaps in software packages, thereby providing users with an accessible method of fitting log-linear models to genotype data on triads.

In \cref{chap:background}, we introduced the genetic background required for modelling genotype data on trios.
We described how the categories corresponding to maternal, paternal, and child genotype values could be modelled with a multinomial distribution with hyperparameters corresponding to genetic and environmental effects.
This also motivated the log-linear model, which allows the probabilities of the multinomial distribution to be captured by a Poisson generalized linear model.
We then discussed the genetic factors that could contribute to disease risk, such as maternal effects, imprinting, and gene-environment interactions, and we showed how they are modelled.

In \cref{chap:methods}, we outlined the features of existing software packages for multinomial (\textbf{EMIM}) and log-linear (\textbf{Haplin}) modelling.
We provided details on the use of \textbf{TriLLIEM} to simulate and model triad data.
Using simulated data with various effect sizes, we compared estimates between \textbf{EMIM} and \textbf{Haplin} to ensure agreement when it was expected.
The differences in functionality between \textbf{EMIM}, \textbf{Haplin}, and \textbf{TriLLIEM} are summarized in \cref{tab:features}.

\begin{table}[htb]
\centering
\caption{Comparison of features between \textbf{Haplin}, \textbf{EMIM}, and \textbf{TriLLIEM}.}
\label{tab:features}
\begin{tblr}{
  colspec={Q[3.1cm] Q[3.7cm] Q[3.7cm] Q[3.7cm]},
  row{1} = {font=\bfseries},
  rowsep=2pt,
  column{1} = {font=\bfseries},
}
\toprule
Functionality & Haplin & EMIM & TriLLIEM \\
\midrule
Availability &
R package, available on CRAN &
Standalone FORTRAN package, compiled from source or Windows/Linux binaries &
R package, available via GitHub \\

Underlying probability model &
Poisson log-linear model &
Multinomial model &
Poisson log-linear model \\

\raggedright Mating type models &
HWE only &
HWE, MS, MaS &
HWE, MS, MaS \\

\raggedright Main genetic effects & \vspace{-1em}
\begin{itemize}[leftmargin=*, noitemsep, topsep=0pt]
    \item Maternal effects
    \item Child effects
\end{itemize} & \vspace{-1em}
\begin{itemize}[leftmargin=*, noitemsep, topsep=0pt]
    \item Maternal effects
    \item Child effects
\end{itemize} & \vspace{-1em}
\begin{itemize}[leftmargin=*, noitemsep, topsep=0pt]
    \item Maternal effects
    \item Child effects
\end{itemize} \\

Parent-of-origin effects & \vspace{-1em}
\begin{itemize}[leftmargin=*, noitemsep, topsep=0pt]
    \item Parent-of-origin effects (modified parameterization)
\end{itemize} & \vspace{-1em}
\begin{itemize}[leftmargin=*, noitemsep, topsep=0pt]
    \item Maternal imprinting effect
    \item Paternal imprinting effect
\end{itemize} & \vspace{-1em}
\begin{itemize}[leftmargin=*, noitemsep, topsep=0pt]
    \item Maternal imprinting effect
    \item Paternal imprinting effect
\end{itemize} \\

Gene-environment interactions &
Stratified approach implemented in the software &
Not implemented, user may code the stratified approach themselves &
Both stratified and non-stratified approaches are implemented \\
\bottomrule
\end{tblr}
\end{table}

The main analysis in this thesis was in \cref{chap:simstudy}, where we conducted two simulation studies using \textbf{TriLLIEM}.
The first was to examine the log-linear model's performance when modelling single populations simulated under many different scenarios (e.g.\ mating asymmetry, true genetic effects), and to compare the stratified and non-stratified approaches to environmental interactions.
The second was to examine the log-linear model's Type 1 error when modelling gene-environment interactions when population stratification is present.
These specific simulations were chosen as they act as extensions to the findings from other researchers.
For instance, imprinting effects were considered in the original Weinberg model \cite{Wein+98}, and further studied in Ainsworth et al.\ \cite{Ains+11}, but work on imprinting by environment interactions is limited to the rigid parameterization used in \textbf{Haplin} \cite{Gjer+18}.
A goal we had with this work was to exploit the flexibility of generalized linear models to explore gene-environment interactions.

From our first simulation study, we confirmed that mating asymmetry results in a biased maternal effect estimator when the HWE or MS mating type model is erroneously assumed.
Furthermore, we observed spurious significance for other parameters (e.g.\ child effects when modelling maternal imprinting by environment interactions) when control triads were present and the data were simulated under mating asymmetry.
We also showed that the stratified approach results in a substantial loss of power for imprinting by environment interactions compared to the non-stratified approach.

From our second simulation study, we found that Type 1 error rates for the main genetic effects were inflated under all mating type models when data were simulated under population stratification with both case- and control-triads when estimating gene-environment interactions.
These interaction terms, however, had Type 1 error rates in the range of $5.5\%$ to $6.7\%$, which is only slightly inflated from the desired $5\%$ rate, and in some cases is within simulation error.
We also found that when the data included only case triads, population stratification did not affect the Type 1 error rates of any of the modelled genetic factors.
These findings corroborated those found by Shin et al.\ \cite{Shin+10}.

There are several future goals we have in mind for improving \textbf{TriLLIEM}.
First, we will allow users to include data with missing counts, and use the EM algorithm to estimate the missing counts to perform the appropriate significance tests, as is possible in \textbf{Haplin} \cite{GjesLie06}.
For instance, there could be genotype failure during data collection that results in the father's genotype being unknown.
Since this is already a feature in \textbf{Haplin} and \textbf{EMIM}, allowing sporadic missing genotype data in \textbf{TriLLIEM} will increase its utility.
We will also implement significance tests beyond the Wald test for individual parameters in the model, such as the likelihood ratio test for comparing nested models.
Although the likelihood is provided by the \texttt{glm()} function, which is used for estimation of genetic effects, when the EM algorithm is required (imprinting, missing genotype data), the likelihood returned by the \texttt{glm()} function in \textsf{R} is not correct.
Several optimizations could be made in \textbf{TriLLIEM}, such as its implementation of the EM algorithm, which runs significantly slower than \textbf{Haplin}'s, and pales in speed compared to \textbf{EMIM}'s direct search.
This can be improved by refactoring some of the relevant code to use a compiled C++ library rather than relying purely on R.
Lastly, we hope to add more features that make \textbf{TriLLIEM} an effective option for users from non-statistical backgrounds, such as functions that automate model selection or provide clear interpretations of the statistical tests and results that \textbf{TriLLIEM} outputs.

\vfill
\cleardoublepage

%%% Local Variables: 
%%% mode: latex
%%% TeX-master: "template"
%%% End: 
